\documentclass[11pt]{article}
\usepackage{amsmath}
\usepackage{hyperref}
\title{Progressive Retrieval Attention with Hierarchical Chunk Routing:\\
A Draft Architecture for Selective Neural Memory}
\author{Draft}
\date{\today}
\begin{document}
\maketitle

\begin{abstract}
This paper presents a draft description of Progressive Retrieval Attention (PRA),
a hierarchical neural-memory architecture in which layer-specific queries
select references, chunks, and cached neural memory progressively rather than
attending to all available context. The paper is intentionally a near-final
draft. Experimental sections contain placeholders pending implementation.
\end{abstract}

\section{Introduction}
Current long-context systems scale computation primarily with available context.
PRA instead aims to scale computation with selected neural memory.

\section{Architecture}
Hierarchy:
\begin{enumerate}
\item Reference selection
\item Chunk selection
\item One routing gist per chunk
\item Materialization of selected chunk token K/V
\item Multi-head memory attention
\end{enumerate}

Routing gists are derived from contextual projected keys by default. Optional
summary metadata may be used through configuration but is never required.

\section{Chunking}
Support:
\begin{itemize}
\item Whole reference
\item Fixed token chunking
\item Marker-based chunking
\item Semantic chunking (future work)
\end{itemize}

\section{Recursive References}
References may recursively resolve other references under bounded depth,
cycle detection and resource budgets.

\section{Dynamic Memory Batching}
Support three modes:
\begin{itemize}
\item Per-example execution
\item Single padded rectangle
\item Dynamic bucketed rectangles grouped by similar memory sizes
\end{itemize}

\section{Training and Retrieval Metrics}
Evaluate language modelling, routing quality, memory utilization,
latency, bucket efficiency and cache reuse independently.

\section{Experimental Plan}
Skeleton:
\begin{itemize}
\item Full-context baseline
\item PRA with full-reference materialization
\item Selected chunk materialization
\item Mean vs last vs GRU routing gists
\item Search strategy comparison
\item Bucket count comparison
\item Recursive reference experiments
\end{itemize}

Results tables and figures intentionally omitted until implementation.

\section{Discussion}
The central hypothesis is that performance can remain close to a
full-context baseline while only a small fraction of available neural memory
is materialized.

\section{Tentative Conclusions}
The proposed architecture unifies hierarchical retrieval, reusable neural
memory, chunk-level routing, recursive references and dynamic execution.
At present this remains a research hypothesis. Whether PRA delivers substantial
quality/computation improvements depends on future empirical evaluation.
The implementation roadmap has been designed to make these questions directly
measurable.

\bibliographystyle{plain}
\begin{thebibliography}{9}
\bibitem{transformer} Vaswani et al. Attention Is All You Need.
\bibitem{retro} RETRO.
\bibitem{rag} Retrieval-Augmented Generation.
\end{thebibliography}
\end{document}
